% Source: https://artofproblemsolving.com/community/c3h2359408_geo_problem_that_gives_you_iq
% Stable identity: aops-topic-2359408-q01
\textbf{Two right triangles.} Since $3^2+4^2=5^2$, triangle $ABC$ has a right angle at $C$, and triangle $ABD$ has a right angle at $D$. Let $E$ be the foot of the perpendicular from $C$ to $AB$. In the right triangle $ABC$, $CE=\dfrac{3\cdot4}{5}=\dfrac{12}{5}$ and $AE=\dfrac{AC^2}{AB}=\dfrac95$. Triangle $ABD$ is congruent to $ABC$ with the same placement of the sides $3$ and $4$ (as $AD=AC=3$ and $BD=BC=4$), so the foot of the perpendicular from $D$ to $AB$ is the same point $E$ and $DE=\dfrac{12}{5}$.

\textbf{A right angle at $E$.} Now $CE^2+DE^2=2\cdot\dfrac{144}{25}=\dfrac{288}{25}=CD^2$, so $\angle CED=90^\circ$.

\begin{asy}
usepackage("amsmath");
size(5cm);
defaultpen(fontsize(10pt));
pair p(real x,real y,real z){return (x+0.6*y,z+0.45*y);}
real e=9/5, k=12/5;
pair A=p(0,0,0), B=p(5,0,0), Ep=p(e,0,0), C=p(e,-k,0), D=p(e,0,k);
draw(A--B,linewidth(0.9)); draw(A--C--B,linewidth(0.9)); draw(A--D--B,linewidth(0.9)); draw(C--D,linewidth(0.9));
draw(C--Ep,dashed+blue); draw(D--Ep,dashed+blue);
dot(A); dot(B); dot(C); dot(D); dot(Ep,blue);
label("$A$",A,W); label("$B$",B,E); label("$C$",C,S); label("$D$",D,N); label("$E$",Ep,SE,blue);
label("$AB=5$",(3.6,0),N,fontsize(8pt)); label("$\dfrac{12}{5}$",(D+Ep)/2,E,blue+fontsize(8pt));
void tick(pair P,pair Q){pair M=(P+Q)/2, v=0.09*unit(rotate(90)*(Q-P)); draw((M-v)--(M+v),blue+linewidth(0.8));}
tick(C,Ep); tick(D,Ep);
pair u1=0.22*unit(D-Ep), u2=0.22*unit(C-Ep), u3=0.22*unit(B-Ep); draw((Ep+u1)--(Ep+u1+u2)--(Ep+u2),blue); draw((Ep+u1)--(Ep+u1+u3)--(Ep+u3),blue);
\end{asy}

\textbf{Volume.} The line $AB$ is perpendicular to both $EC$ and $ED$, hence to the plane $CED$. So the tetrahedron is made of two pyramids with base $CED$ and heights $AE$ and $EB$, which add up to $AB=5$:
\[V=\dfrac13\cdot[CED]\cdot AB=\dfrac13\cdot\dfrac12\cdot\dfrac{12}{5}\cdot\dfrac{12}{5}\cdot5=\boxed{\dfrac{24}{5}}.\]
